% Source: sources/returns/S6_04_MAPS_MANIFOLDS/sections/manifolds.tex:8492–8802
\paragraph{Theorem D.2.43 (Proved: a general triangular cusp and its full boundary kernel).}
Let $N_L$ be the smooth same-fan triangular Mumford cusp
whose short-period lattice is a full sublattice
$L\subset\mathbb Z^2$. Write its Smith factors as
$s\mid t$, and put $D=st$. Holomorphic unit factors
in the two periods are allowed. Let $M_L=\partial N_L$,
and let $\Lambda=\mathbb Z^4$ be the regular torus
lattice. Its unipotent monodromy is $A=I+N$,
where $N^2=0$ and the map from the long lattice to
the primitive short lattice has Smith factors $(s,t)$.
Put $F=\ker(A-I)$, the primitive short lattice.

Then $\pi_1(N_L)=L\simeq\mathbb Z^2$, and its ordinary
integral homology is
\[
 \begin{array}{c|ccccc}
 j&0&1&2&3&4\\ \hline
 H_j(N_L)&\mathbb Z&\mathbb Z^2&
       \mathbb Z^{D+3}\oplus\mathbb Z/s&
       \mathbb Z^2&\mathbb Z^D .
 \end{array}                                             \tag{D.GC1}
\]
All higher groups vanish. In particular $\chi(N_L)=2D$.
The full regular-fibre kernel is the generally
non-saturated lattice
\[
 \ker\{H_2(T^4)\to H_2(N_L)\}
       =(\Lambda^2A-I)\Lambda^2\Lambda+\Lambda^2F .
                                                        \tag{D.GC2}
\]
Its image is $\mathbb Z^4\oplus\mathbb Z/s$, with
free cokernel of rank $D-1$ in $H_2(N_L)$.
The two actual fixed angular normal tori bound
and give the entire fixed-normal projection.
Together with (D.GC2) they specify the full
boundary kernel. More explicitly,
\[
 \begin{split}
 H_1(M_L)&=\mathbb Z^3\oplus\mathbb Z/s\oplus\mathbb Z/t,\\
 H_2(M_L)&=\mathbb Z^6\oplus(\mathbb Z/s)^2,\\
 \ker\{H_2(M_L)\to H_2(N_L)\}
             &\simeq\mathbb Z^2\oplus\mathbb Z/s .
 \end{split}                                             \tag{D.GC3}
\]
The torsion killed by this map is the primitive
angular area modulo the monodromy image. The other
boundary torsion summand maps isomorphically to
the actual cap torsion.

There is also an exact component-meridian marking.
Use the positive $q$-meridian $\theta$ and index
the $D$ normalized components by $v\in\mathbb Z^2/L$.
The positive transverse disc to that component has
boundary
\[
                     \theta+\bar v
           \quad\hbox{in }\mathbb Z\theta\oplus\mathbb Z^2/L
                     \subset H_1(M_L).                  \tag{D.GC4}
\]
These are the complete columns of the relative
map $H_2(N_L,M_L)=H^4(N_L)=\mathbb Z^D\to H_1(M_L)$.
In particular the component differences need not
vanish before global gluing.

\begin{proof}
\emph{The actual covering and its two-variable normalization.}
The extension of period units and the free same-fan
cover construction in D.2.22 identify $N_L$ with
the finite cover of the minimal hexagonal cusp
associated to $L$. The units can also be removed
through the relative smooth family of E.2.19.
The normalized components are $dP_6$ surfaces,
the double curves are spheres, and the triple
strata are points. Their dual complex is
$\mathbb R^2/L$. Van Kampen gives
$\pi_1(N_L)=L$, because each component and every
intersection component is simply connected.

View the common universal cover with the full deck
action of the minimal cusp, and put
$R=\mathbb Z[x^{\pm1},y^{\pm1}]$ and
$I=(x-1,y-1)$. In the cyclic toric order
$(E_1,H-E_1-E_2,E_2,H-E_2-E_3,E_3,H-E_1-E_3)$,
the three opposite-pair translations have
weights $x^{-1},y,xy$. Indeed the primitive
fan rays are $(1,0),(0,1),(-1,1)$, and the
opposite pairing translates by their negatives;
choose the second period generator with the
corresponding sign. Simultaneously inverting
the deck variables changes neither the ideal
$I^2$ nor the assertions below.

With intersection form $\operatorname{diag}(1,-1,-1,-1)$,
the degree-two normalization matrix is
\[
 B=\begin{pmatrix}
 -x^{-1}&-1&-x^{-1}&-x^{-1}\\
 1&1&1&y\\
 -xy&-xy&-1&-xy
 \end{pmatrix}.                                         \tag{D.GC5}
\]
The homology differential is its transpose with
the corresponding deck involution. One unit
row/column cancellation reduces its transpose to
\[
 F=\begin{pmatrix}
 1-x&0\\0&xy-1\\y-1&xy(y-1)
 \end{pmatrix}.                                         \tag{D.GC6}
\]
For clarity there is an explicit Laurent-unimodular
comparison, not just equality of ranks or minors.
Let
\[
 T=\begin{pmatrix}
 0&0&-x^{-1}\\
 -x^{-1}y^{-2}&x^{-1}y^{-1}&-x^{-2}y^{-1}\\
 -y^{-1}&0&0
 \end{pmatrix},\qquad
 S=\begin{pmatrix}y&xy^2\\x&0\end{pmatrix}.
\]
Their determinants are Laurent units, and
\[
 T^{-1}F=
 \begin{pmatrix}-(y-1)&0\\x-1&-(y-1)\\0&x-1\end{pmatrix}S.
                                                        \tag{D.GC7}
\]
The middle matrix is the exact presentation
of $I^2$ by the generators
$(x-1)^2,(x-1)(y-1),(y-1)^2$.
Its exactness follows directly from the regular
sequence $x-1,y-1$: the two displayed relations
are all the syzygies among those three monomials.
Thus
\[
                H_2(\widetilde N_L;\mathbb Z)\simeq I^2
                         \quad\hbox{as an }R\hbox{-module}.
                                                        \tag{D.GC8}
\]
Here the universal cover is the same for every $L$.
The degree-zero normalization row is the cellular
complex of the plane, and the degree-four row is
$R$ in normalization degree zero. The degree-two
differential is injective over the domain $R$.
No higher differential is possible, because
normalization degree is at most two and the odd
homology of every stratum vanishes.

\emph{The integral finite-cover groups.}
Write $G=\mathbb Z^2/L$ and $R_G=\mathbb Z[G]$.
The degree-two normalization cokernel is
$I^2\otimes_R R_G$. We compute its torsion exactly.
The two short exact sequences
\[
 0\to I\to R\to\mathbb Z\to0,\qquad
 0\to I^2\to I\to\mathbb Z^2\to0
                                                        \tag{D.GC9}
\]
use the trivial $R$-action on $I/I^2=\mathbb Z^2$.
The Koszul resolution for $x-1,y-1$ identifies
$\operatorname{Tor}^R_j(\mathbb Z,R_G)$ with
$H_j(L;\mathbb Z)$.
Consequently $I\otimes_R R_G$ is free of rank $D+1$:
it is an extension of the free augmentation ideal
of $\mathbb Z[G]$, of rank $D-1$, by $L\simeq\mathbb Z^2$.
Its map to $\mathbb Z^2$ in (D.GC9) is onto,
since $x-1$ and $y-1$ map to the two standard
generators. Its kernel is free of rank $D-1$.

The relevant Tor segment from the second sequence is
\[
 \mathbb Z\ \xrightarrow{\;\alpha_L\;}\ \mathbb Z^4
 \longrightarrow I^2\otimes_RR_G
 \longrightarrow I\otimes_RR_G
 \longrightarrow\mathbb Z^2\longrightarrow0 .
                                                        \tag{D.GC10}
\]
The source is $H_2(L)$ and the target is
$H_1(L)\otimes\mathbb Z^2$.
If $u,v$ are an oriented basis of $L$, the map is
\[
                  \alpha_L(1)=u\otimes v-v\otimes u,
                                                        \tag{D.GC11}
\]
where the second tensor factor uses the inclusion
$L\hookrightarrow\mathbb Z^2$.
One can check (D.GC11) directly in the Koszul
complex. In Smith coordinates
$L=\langle(s,0),(0,t)\rangle$, let
$N_x=1+\cdots+x^{s-1}$ and
$N_y=1+\cdots+y^{t-1}$.
The map on resolutions sends the relation
$(-(y-1),x-1)$ of $I$ to
$-e_y\otimes e_x+e_x\otimes e_y$.
The invariant cycle $N_xN_y$ therefore maps,
in $H_1(L)\otimes\mathbb Z^2$, to a vector with
nonzero coordinates $-s,t$.
Its Smith factor is exactly $s=\gcd(s,t)$.
The augmentation ideal and its square are invariant
under changes of basis of $\mathbb Z^2$, so this
also proves (D.GC11) for an arbitrary $L$.

Thus $\alpha_L$ is injective and its cokernel is
$\mathbb Z^3\oplus\mathbb Z/s$.
The free last kernel in (D.GC10) gives
\[
            I^2\otimes_RR_G
                    \simeq\mathbb Z^{D+2}\oplus\mathbb Z/s.
                                                        \tag{D.GC12}
\]
The preceding Tor segment gives
$\operatorname{Tor}_1^R(I^2,R_G)=\mathbb Z^2$.
The degree-zero normalization row now has the
homology of the dual torus, adding one free
summand to total degree two and two to degree
one. The degree-four row contributes $\mathbb Z^D$.
These observations prove (D.GC1); the extension
by the dual torus's degree-two $\mathbb Z$ splits
as an abelian-group extension.

\emph{The actual boundary map and the nonsaturated kernel.}
Choose bases with
$A b_1=b_1+s a_1$, $A b_2=b_2+t a_2$ and
$A a_i=a_i$. The exterior-square difference has
image generated by
\[
 s\,a_1\wedge a_2,\qquad
 s\,a_1\wedge b_2-t\,a_2\wedge b_1
                         +st\,a_1\wedge a_2 .
                                                        \tag{D.GC13}
\]
Its two Smith factors are $(s,s)$.
The Wang sequence gives the first two boundary
groups in (D.GC3).

The minimal cusp's two fixed angular normal
tori have actual section-disc bounding sweeps.
In the unitless toric chart at the dense locus
of the component indexed by zero, take the
section with its two nonvanishing torus
coordinates equal to one and its normal
coordinate equal to $q$. Acting on this
section by either short circle gives a map
$S^1\times D^2\to N$ whose boundary is the
literal fixed angular normal torus. Its
fundamental circle lies in a simply connected
normalized component, so its image in the
cusp fundamental group is zero. The whole
map therefore lifts to every period cover
$N_L$, with its marked boundary. Transport
through the relative unit-removing family
preserves these boundary classes. Thus the
two primitive fixed-normal projections still
bound, and their actual classes split off
the full Wang projection $H_2(M_L)\to F$.

The primitive angular area also maps to zero.
In the universal cover the lifted regular
fibre is $T^2\times\mathbb R^2$. Its degree-two
group is a trivial $R$-module. Its image in
the module $I^2$ of (D.GC8) is zero, since
$I^2$ has no nonzero element fixed by both
$x$ and $y$. This nullhomology projects to
every finite period cover.

Poincare--Lefschetz duality gives
\[
 H_3(N_L,M_L)=H^3(N_L)=\mathbb Z^2\oplus\mathbb Z/s,
 \qquad H_2(N_L,M_L)=H^4(N_L)=\mathbb Z^D .
                                                        \tag{D.GC14}
\]
The already bounding angular normal tori exhaust
the rational rank of the boundary kernel.
Hence the four-dimensional rational fibre
coinvariant injects into the cap.
Moreover the cokernel of
$H_2(M_L)\to H_2(N_L)$ injects into the free
group $H^4(N_L)$, so it is free.
Because the fixed-normal projections bound,
the boundary image equals the regular-fibre image.
It therefore contains all of the cap torsion.
There is no hidden free-to-torsion extension
here: if a fibre-coinvariant element maps to
a torsion class, rational injectivity forces
that element itself to be torsion. Hence
the cap torsion is the image of the boundary
fibre torsion subgroup, not of a free class
whose rational image has been overlooked.

The two boundary fibre torsion summands in
(D.GC13) have order $s$. The angular-area
one is killed. The other must map onto
$\operatorname{Tor}H_2(N_L)=\mathbb Z/s$ and
hence maps isomorphically. No free fibre
class can enter the kernel, since the rational
fibre coinvariant injects. This proves exactly
(D.GC2), not its saturation, and gives all
the remaining assertions in (D.GC3).
The free cokernel has rank $(D+3)-4=D-1$.

Finally the smooth fan has primitive rays
$(v,1)$. A positive small normal disc to
the associated component has this primitive
angular-plus-base meridian. Reducing the
angular coordinate modulo the period lattice
gives (D.GC4), with the complex orientation
fixing the common sign. Every coset of $L$
occurs once. Their differences generate
$\mathbb Z^2/L$, and the zero coset gives
$\theta$; thus the displayed columns give
the entire kernel of $H_1(M_L)\to H_1(N_L)$.
\end{proof}
